\chapter{APÊNDICE D — DADOS DE ENTRADA E MEMÓRIA DE CÁLCULO DAS ALTURAS E INCLINAÇÕES}\label{ap:dados-morfometria}

Este apêndice reúne as coordenadas, os tempos duplos de percurso e os parâmetros utilizados no cálculo dos desníveis entre A e C e das inclinações do declive não deformado e da superfície receptora. As medidas foram obtidas diretamente na interpretação sísmica, considerando o tempo do horizonte e o tempo do fundo marinho na mesma posição. A localização dos pontos e dos segmentos é apresentada nas figuras dos respectivos depósitos no capítulo de resultados.

As denominações anteriores correspondem a: Santos — MTD 1; Continuação Santos — MTD 2; Recorrência — MTD 3; Regional recente — evento C do MTC 1; e Búzios recente — evento D do MTC 2.

\section{Identificação e localização dos pontos}

A e C correspondem aos extremos utilizados no cálculo do desnível longitudinal. A identificação de A como topo da cicatriz ou extremo proximal do depósito deve acompanhar a interpretação de cada caso. C corresponde ao ponto basal na terminação distal. No evento D do MTC 2, os dois eixos utilizam o mesmo ponto A, correspondente ao ponto A originalmente definido para o eixo 1. Apenas as terminações distais diferem, sendo identificadas como C1 e C2.

P1 e P2 delimitam o segmento do declive não deformado, externo à zona deformada. P3 e P4 delimitam o segmento da superfície receptora, à frente do pé do depósito.

\begin{table}[htbp]
  \centering
  \caption{Coordenadas dos extremos utilizados nos cálculos entre A e C.}
  \label{tab:resultado-D-1}
  \scriptsize
  \setlength{\tabcolsep}{2.5pt}
  \begin{tabularx}{\textwidth}{>{\raggedright\arraybackslash}X >{\centering\arraybackslash}X >{\centering\arraybackslash}X >{\centering\arraybackslash}X >{\centering\arraybackslash}X >{\centering\arraybackslash}X}
    \toprule
    \textbf{Depósito/evento} & \textbf{Trajetória} & \textbf{X de A (m)} & \textbf{Y de A (m)} & \textbf{X de C (m)} & \textbf{Y de C (m)} \\
    \midrule
    MTD 1 & A--C & 778.211,95 & 7.325.273,75 & 777.578,48 & 7.319.161,82 \\
    MTD 2 & A--C & 761.062,25 & 7.316.233,28 & 761.586,26 & 7.313.836,80 \\
    MTD 3 & A--C & 703.403,00 & 7.280.395,99 & 699.993,75 & 7.265.391,56 \\
    MTC 2-D & Eixo 1 — A comum--C1 & 740.770,81 & 7.289.216,64 & 742.886,64 & 7.264.223,92 \\
    MTC 2-D & Eixo 2 — A comum--C2 & 740.770,81 & 7.289.216,64 & 752.202,64 & 7.275.546,87 \\
    \bottomrule
  \end{tabularx}
\end{table}

\nota{no MTC 2-D, adotou-se provisoriamente o ponto inicial do eixo 1 como origem comum das duas trajetórias. As coordenadas são apresentadas conforme os dados fornecidos. O sistema de referência de coordenadas deve ser identificado na versão final.}

\begin{table}[htbp]
  \centering
  \caption{Coordenadas dos pontos utilizados nas inclinações do declive não deformado e da superfície receptora.}
  \label{tab:resultado-D-2}
  \scriptsize
  \setlength{\tabcolsep}{2.5pt}
  \begin{tabularx}{\textwidth}{>{\raggedright\arraybackslash}X >{\centering\arraybackslash}X >{\centering\arraybackslash}X >{\centering\arraybackslash}X >{\centering\arraybackslash}X >{\centering\arraybackslash}X}
    \toprule
    \textbf{Depósito/evento} & \textbf{Segmento} & \textbf{X inicial (m)} & \textbf{Y inicial (m)} & \textbf{X final (m)} & \textbf{Y final (m)} \\
    \midrule
    MTD 1 & P1--P2 — S & 769.508,00 & 7.323.203,00 & 770.451,00 & 7.317.402,00 \\
    MTD 1 & P3--P4 — St & 777.605,63 & 7.318.494,15 & 776.640,00 & 7.318.485,85 \\
    MTD 2 & P1--P2 — S & 756.961,40 & 7.314.907,08 & 757.822,91 & 7.312.473,92 \\
    MTD 2 & P3--P4 — St & 761.217,68 & 7.313.478,47 & 760.815,19 & 7.312.462,46 \\
    MTD 3 & P1--P2 — S & 699.105,50 & 7.279.766,38 & 697.664,31 & 7.266.102,66 \\
    MTC 1-C & P1--P2 — S & 740.348,72 & 7.336.160,47 & 752.724,58 & 7.305.074,70 \\
    MTC 1-C & P3--P4 — St & 735.992,35 & 7.294.585,64 & 737.582,27 & 7.291.812,31 \\
    \bottomrule
  \end{tabularx}
\end{table}

\nota{inicial e final correspondem, respectivamente, a P1 e P2 para S e a P3 e P4 para St.}

\section{Tempos utilizados nos cálculos}

As Tabelas D.3 e D.4 apresentam os tempos individuais dos pontos interpretados e do fundo marinho nas mesmas posições. Todos os valores são expressos em milissegundos de tempo duplo de percurso — ms TWT.

\begin{table}[htbp]
  \centering
  \caption{Tempos dos extremos utilizados nos cálculos entre A e C.}
  \label{tab:resultado-D-3}
  \scriptsize
  \setlength{\tabcolsep}{2.5pt}
  \begin{tabularx}{\textwidth}{>{\raggedright\arraybackslash}X >{\centering\arraybackslash}X >{\centering\arraybackslash}X >{\centering\arraybackslash}X >{\centering\arraybackslash}X >{\centering\arraybackslash}X}
    \toprule
    \textbf{Depósito/evento} & \textbf{Trajetória} & \textbf{TWT de A (ms)} & \textbf{Fundo em A (ms)} & \textbf{TWT de C (ms)} & \textbf{Fundo em C (ms)} \\
    \midrule
    MTD 1 & A--C & 1.872,44 & 1.819,73 & 2.190,13 & 2.079,84 \\
    MTD 2 & A--C & 1.979,96 & 1.938,34 & 2.147,95 & 2.059,03 \\
    MTD 3 & A--C & 2.358,80 & 2.139,12 & 2.897,46 & 2.561,82 \\
    MTC 2-D & Eixo 1 — A comum--C1 & 2.179,34 & 2.126,42 & 2.926,58 & 2.808,03 \\
    MTC 2-D & Eixo 2 — A comum--C2 & 2.179,34 & 2.126,42 & 2.774,16 & 2.556,13 \\
    \bottomrule
  \end{tabularx}
\end{table}

\nota{os tempos de A são compartilhados pelos dois eixos do MTC 2-D.}

\begin{table}[htbp]
  \centering
  \caption{Tempos utilizados nos cálculos das inclinações.}
  \label{tab:resultado-D-4}
  \scriptsize
  \setlength{\tabcolsep}{2.5pt}
  \begin{tabularx}{\textwidth}{>{\raggedright\arraybackslash}X >{\centering\arraybackslash}X >{\centering\arraybackslash}X >{\centering\arraybackslash}X >{\centering\arraybackslash}X >{\centering\arraybackslash}X}
    \toprule
    \textbf{Depósito/evento} & \textbf{Segmento} & \textbf{TWT inicial do horizonte (ms)} & \textbf{Fundo no início (ms)} & \textbf{TWT final do horizonte (ms)} & \textbf{Fundo no final (ms)} \\
    \midrule
    MTD 1 & P1--P2 — S & 2.116,71 & 2.032,27 & 2.199,62 & 2.083,20 \\
    MTD 1 & P3--P4 — St & 2.183,64 & 2.070,47 & 2.188,42 & 2.078,42 \\
    MTD 2 & P1--P2 — S & 2.003,88 & 1.910,67 & 2.224,98 & 2.127,43 \\
    MTD 2 & P3--P4 — St & 2.140,21 & 2.056,21 & 2.133,18 & 2.059,09 \\
    MTD 3 & P1--P2 — S & 2.367,66 & 2.096,71 & 2.738,26 & 2.455,11 \\
    MTC 1-C & P1--P2 — S & 1.194,48 & 1.125,68 & 2.269,85 & 2.047,60 \\
    MTC 1-C & P3--P4 — St & 2.236,91 & 1.978,65 & 2.249,61 & 2.029,46 \\
    \bottomrule
  \end{tabularx}
\end{table}

\section{Parâmetros de conversão e distâncias}

A velocidade da água foi mantida fixa em 1.500 m/s. Para o percurso sedimentar, foram utilizadas as velocidades representativas de cada evento e suas dispersões ponderadas, conforme a metodologia.

\begin{table}[htbp]
  \centering
  \caption{Velocidades sedimentares e componentes de incerteza utilizadas na conversão.}
  \label{tab:resultado-D-5}
  \scriptsize
  \setlength{\tabcolsep}{2.5pt}
  \begin{tabularx}{\textwidth}{>{\raggedright\arraybackslash}X >{\centering\arraybackslash}X >{\centering\arraybackslash}X}
    \toprule
    \textbf{Depósito/evento} & \textbf{Velocidade sedimentar (m/s)} & \textbf{Dispersão ponderada — uv (m/s)} \\
    \midrule
    MTD 1 & 1.561,76 & 13,57 \\
    MTD 2 & 1.554,90 & 18,12 \\
    MTD 3 & 1.715,30 & 32,90 \\
    MTC 1-C & 1.552,59 & 18,86 \\
    MTC 2-D — ambos os eixos & 1.561,17 & 14,15 \\
    \bottomrule
  \end{tabularx}
\end{table}

\nota{valores de entrada apresentados com as casas disponíveis, sem aplicação do arredondamento dos resultados finais.}

\begin{table}[htbp]
  \centering
  \caption{Distâncias horizontais utilizadas nas inclinações e comprimentos dos eixos do MTC 2-D.}
  \label{tab:resultado-D-6}
  \scriptsize
  \setlength{\tabcolsep}{2.5pt}
  \begin{tabularx}{\textwidth}{>{\raggedright\arraybackslash}X >{\centering\arraybackslash}X >{\centering\arraybackslash}X}
    \toprule
    \textbf{Depósito/evento} & \textbf{Medida} & \textbf{Distância ou comprimento (km)} \\
    \midrule
    MTD 1 & Segmento P1--P2 & 5,877 \\
    MTD 1 & Segmento P3--P4 & 0,966 \\
    MTD 2 & Segmento P1--P2 & 2,581 \\
    MTD 2 & Segmento P3--P4 & 1,093 \\
    MTD 3 & Segmento P1--P2 & 13,740 \\
    MTC 1-C & Segmento P1--P2 & 33,459 \\
    MTC 1-C & Segmento P3--P4 & 3,197 \\
    MTC 2-D & Comprimento total do eixo 1 & 24,33 \\
    MTC 2-D & Comprimento total do eixo 2 & 23,48 \\
    \bottomrule
  \end{tabularx}
\end{table}

\nota{as distâncias dos segmentos foram mantidas conforme as medidas fornecidas. Cada eixo do MTC 2-D inclui 9 km de trecho comum e sua respectiva ramificação. As distâncias P1--P2 e P3--P4 não correspondem aos comprimentos totais dos eixos.}
